% ════════════════════════════════════════════════════════════════
\section{Programowanie dynamiczne}
% ════════════════════════════════════════════════════════════════

Naiwna implementacja rekurencyjna bywa wykładnicza, bo te same podproblemy liczone są wielokrotnie
(np. \Call{Fib}{$n$} przelicza $F_{n-1}, F_{n-2}, \dots$ od zera). Programowanie dynamiczne rozwiązuje
każdy podproblem raz i zapamiętuje wynik w tablicy, schodząc z czasu wykładniczego do wielomianowego.
Stosuje się je, gdy problem ma \emph{optymalną podstrukturę} (rozwiązanie optymalne składa się
z~optymalnych rozwiązań podproblemów) i~\emph{nakładające się podproblemy}.

% ════════════════════════════════════════════════════════════════
\subsection{Algorytm mnożenia macierzy}
% ════════════════════════════════════════════════════════════════

Przypomnijmy, jak mnoży się macierze. Iloczyn $C = A \cdot B$ macierzy $A$ rozmiaru $p \times q$
oraz $B$ rozmiaru $q \times r$ jest macierzą $C$ rozmiaru $p \times r$, w~której element $C(i, j)$
to iloczyn skalarny $i$-tego wiersza macierzy $A$ oraz $j$-tej kolumny macierzy $B$:
\[
	C(i, j) = \sum_{k = 1}^{q} A(i, k) \cdot B(k, j).
\]
Biorąc dla przykładu $p = 2$, $q = 3$, $r = 2$ i~obliczając wyróżniony element $C(1, 1)$:
\[
	\underbrace{\begin{bmatrix}
			{\color{accentA}1} & {\color{accentA}2} & {\color{accentA}3} \\
			4                  & 5                  & 6
		\end{bmatrix}}_{A:\ p \times q}
	\cdot
	\underbrace{\begin{bmatrix}
			{\color{accentB}7}  & 8  \\
			{\color{accentB}9}  & 10 \\
			{\color{accentB}11} & 12
		\end{bmatrix}}_{B:\ q \times r}
	=
	\underbrace{\begin{bmatrix}
			{\color{accentA}58} & 64  \\
			139                 & 154
		\end{bmatrix}}_{C:\ p \times r}
\]
\[
	C(1, 1) = {\color{accentA}1} \cdot {\color{accentB}7} + {\color{accentA}2} \cdot {\color{accentB}9}
	+ {\color{accentA}3} \cdot {\color{accentB}11} = 7 + 18 + 33 = 58.
\]
Kolorowy wiersz macierzy $A$ i~ kolorowa kolumna macierzy $B$ dają kolorowy element macierzy $C$;
sumowanie po~$k = 1, \dots, q$ to wewnętrzna pętla algorytmu. Każdy z~$p \cdot r$ elementów macierzy
$C$ wymaga $q$ mnożeń.

\begin{alg}{Naiwne mnożenie macierzy}{alg:matmul}
	\FunctionNN{MatrixMultiply}{$A, B$}
	\Comment{$A$ -- macierz $p \times q$; $B$ -- macierz $q \times r$; $C$ -- macierz $p \times r$}
	\For{$i \Assign 1$ \To $p$}
	\For{$j \Assign 1$ \To $r$}
	\State $C(i, j) \Assign 0$
	\For{$k \Assign 1$ \To $q$}
	\State $C(i, j) \Assign C(i, j) + A(i, k) \cdot B(k, j)$ \Comment{operacja dominująca}
	\EndFor
	\EndFor
	\EndFor
	\State \Return $C$
	\EndFunction
\end{alg}

Operacją dominującą jest mnożenie w linii 5. Mamy $p \cdot q \cdot r$ mnożeń. Złożoność jest $\BO(p \cdot q \cdot r)$.

% ════════════════════════════════════════════════════════════════
\subsection{Problem mnożenia ciągu macierzy}
% ════════════════════════════════════════════════════════════════

\begin{example}[Mnożenie ciągu macierzy]
	$(A_1, A_2, A_3)$ gdzie $A_1$ -- macierz $5 \times 10$, $A_2$ -- macierz $10 \times 3$, $A_3$ -- macierz $3 \times 2$.
	Ile potrzeba mnożeń skalarnych (mnożenia liczb) do obliczenia iloczynu $A_1 \cdot A_2 \cdot A_3$?
	\begin{flalign*}
		A_1 \cdot A_2 \cdot A_3 = (A_1 \cdot A_2) \cdot A_3 = A_1 \cdot (A_2 \cdot A_3) &  &  &  &
	\end{flalign*}

	\begin{itemize}
		\item[I sposób:] $(A_1 \cdot A_2) \cdot A_3 = 5 \cdot 10 \cdot 3 + 5 \cdot 3 \cdot 2 = 150 + 30 = 180$
		\item[II sposób:] $A_1 \cdot (A_2 \cdot A_3) = 10 \cdot 3 \cdot 2 + 5 \cdot 10 \cdot 2 = 60 + 100 = 160$
	\end{itemize}
	II sposób jest szybszy.
\end{example}

\begin{problem}[mnożenia ciągu macierzy]
	Dla danego ciągu macierzy $(A_1, A_2, \dots, A_n)$, gdzie $A_i$ jest macierzą $p_{i-1} \times p_i$ (cały ciąg wymiarów zapisujemy jednym ciągiem $p = (p_0, p_1, \dots, p_n)$ -- kolejne macierze dzielą wspólny wymiar, więc $n$ macierzy opisuje $n + 1$ liczb), znajdź kolejność wykonywania mnożeń minimalizującą liczbę mnożeń skalarnych użytych do obliczenia iloczynu $A_1 \cdot A_2 \cdot \dots \cdot A_n$.
\end{problem}

Ile jest różnych sposobów obliczenia tego iloczynu? Innymi słowy, na ile sposobów możemy poprawnie wstawić nawiasy w iloczynie $A_1 \cdot A_2 \cdot \dots \cdot A_n$?

\begin{example}[Kolejności nawiasowania dla $n=4$]
	$A_1 \cdot A_2 \cdot A_3 \cdot A_4$ \quad ($n = 4$)
	\begin{itemize}
		\item $((A_1 \cdot A_2) \cdot A_3) \cdot A_4$
		\item $(A_1 \cdot (A_2 \cdot A_3)) \cdot A_4$
		\item $A_1 \cdot ((A_2 \cdot A_3) \cdot A_4)$
		\item $A_1 \cdot (A_2 \cdot (A_3 \cdot A_4))$
		\item $(A_1 \cdot A_2) \cdot (A_3 \cdot A_4)$
	\end{itemize}
	Zatem mamy 5 sposobów.
\end{example}

Sprawdźmy czy umiemy rozwiązać ten problem szybko poprzez sprawdzenie wszystkich możliwości.
Niech $P(n)$ -- liczba sposobów wstawienia nawiasów w iloczynie $n$ macierzy.
Załóżmy, że ostatnie mnożenie w ciągu $A_1 \dots A_n$ to mnożenie między $A_k$ i $A_{k+1}$:
$(A_1 \dots A_k)(A_{k+1} \dots A_n)$. W podciągach lewym i prawym nawiasy są rozmieszczane niezależnie od siebie.

Stąd rekurencja:
\begin{flalign*}
	 & \begin{cases}
		   P(n) = \displaystyle\sum_{k=1}^{n-1} P(k) \cdot P(n - k) \\
		   P(1) = 1 \vphantom{\displaystyle\sum_{k=1}^{n-1}}
	   \end{cases} &  &  &  &
\end{flalign*}
Jest to znana formuła rekurencyjna a rozwiązanie to przesunięta liczba Catalana:
\begin{flalign*}
	P(n) = C(n - 1), \text{ gdzie } C(n) = \frac{1}{n + 1}\binom{2n}{n} \text{ jest n-tą liczbą Catalana.} &  &  &  &
\end{flalign*}
Wiadomo, że $C(n) = \BM\left(\frac{4^n}{n^{3/2}}\right)$. Zatem $P(n)$ jest wykładniczy względem $n$.
Chcemy znaleźć szybki algorytm.

\paragraph{Rozwiązanie z optymalną podstrukturą:}
Zauważmy że jeśli optymalne rozwiązanie jest postaci $(A_1 \dots A_k)(A_{k+1} \dots A_n)$ dla jakiegoś $k \in \{1, \dots, n-1\}$, to w podciągach $A_1 \dots A_k$ i $A_{k+1} \dots A_n$ rozwiązanie nawiasów musi być optymalne. Wpp. lepsze rozmieszczenie nawiasów w $A_1 \dots A_k$ lub $A_{k+1} \dots A_n$ dałoby lepsze rozmieszczenie nawiasów w całym ciągu.

Rozwiązanie optymalne zawiera rozwiązania optymalne podproblemów (własność optymalnej podstruktury). Podproblemami w naszym problemie są problemy optymalnego rozmieszczenia nawiasów w ciągach $A_i A_{i+1} \dots A_j$ dla $1 \leq i \leq j \leq n$.

Niech $m(i, j)$ oznacza minimalną liczbę mnożeń skalarnych potrzebnych do obliczenia tego iloczynu $A_i \dots A_j$ dla $1 \leq i \leq j \leq n$. Chcemy obliczyć $m(1, n)$.
\begin{flalign*}
	m(i, i) = 0 &  &  &  &
\end{flalign*}
Optymalny sposób rozmieszczenia nawiasów jest postaci $(A_i \dots A_k)(A_{k+1} \dots A_j)$ dla pewnego $k \in \{i, \dots, j-1\}$. Stąd mamy zależność:
\begin{flalign*}
	m(i, j) = \min_{i \leq k < j} \left\{ m(i, k) + m(k + 1, j) + p_{i-1} \cdot p_k \cdot p_j \right\} \quad \text{dla } i < j &  &  &  &
\end{flalign*}

Niech $s(i, j)$, dla $1 \leq i < j \leq n$, będzie wartością $k$, dla którego to minimum jest osiągane (czyli $s(i, j)$ wskazuje miejsce ostatniego mnożenia w optymalnym rozwiązaniu podproblemu $A_i \dots A_j$).

Jeśli zaimplementujemy po prostu tę formułę rekurencyjną otrzymamy alg. wykładniczy. Ale jest tylko $n + \binom{n}{2} = \BO(n^2)$ podproblemów, więc rzędu $\BO(n^2)$ liczb $m(i, j)$ i każdą z tych liczb wystarczy obliczyć jeden raz.

Pomysł programowania dynamicznego polega więc na zapamiętaniu tych wartości w~tablicy zamiast liczenia ich wielokrotnie. Wprowadzamy dwie tablice indeksowane parą $(i, j)$, $1 \leq i \leq j \leq n$ (a~ponieważ rozważamy tylko $i \leq j$, wypełniona zostaje jedynie górna połowa każdej z~nich -- jak w~przykładzie poniżej):
\begin{itemize}
	\item $m(i, j)$ -- \textbf{koszt}: minimalna liczba mnożeń skalarnych dla podproblemu $A_i \dots A_j$. Ostateczną odpowiedzią jest $m(1, n)$.
	\item $s(i, j)$ -- \textbf{cięcie}: miejsce $k$ ostatniego mnożenia, na którym ten minimalny koszt jest osiągany. Sama tablica $m$ podaje tylko \emph{ile} mnożeń wystarczy, ale nie \emph{jak} ustawić nawiasy; to $s$ pozwala odtworzyć optymalne nawiasowanie (wykorzysta ją algorytm \Call{MatrixChainMultiply}{} dalej).
\end{itemize}

\begin{alg}{Algorytm Matrix Chain Order}{alg:matrixchainorder}
	\FunctionNN{MatrixChainOrder}{$p$}
	\Comment{$p = (p_0, p_1, \dots, p_n)$ - ciąg rozmiarów}
	\State $n \Assign \Length{p} - 1$ \Comment{$A_i$ - macierz $p_{i-1} \times p_i$}
	\For{$i \Assign 1$ \To $n$}
	\State $m(i, i) \Assign 0$
	\EndFor
	\For{$l \Assign 2$ \To $n$} \Comment{$l$ -- liczba macierzy w podproblemie}
	\For{$i \Assign 1$ \To $n - l + 1$} \Comment{$l = 2:\ A_1A_2,\ A_2A_3, \dots, A_{n-1}A_n$}
	\State $j \Assign i + l - 1$ \Comment{$l = 3:\ A_1A_2A_3,\ A_2A_3A_4, \dots, A_{n-2}A_{n-1}A_n$}
	\State $m(i, j) \Assign \infty$ \Comment{$l = n:\ A_1 \dots A_n$}
	\For{$k \Assign i$ \To $j - 1$}
	\State $q \Assign m(i, k) + m(k + 1, j) + p_{i-1} \cdot p_k \cdot p_j$
	\If{$q < m(i, j)$}
	\State $m(i, j) \Assign q$
	\State $s(i, j) \Assign k$
	\EndIf
	\EndFor
	\EndFor
	\EndFor
	\State \Return $m, s$
	\EndFunction
\end{alg}

Mamy 3 pętle (po $l, i, k$) i liczbę iteracji każdej z tych pętli można ograniczyć przez $n$, zatem złożoność czasowa to $\BO(n^3)$.

\begin{example}[Optymalna kolejność mnożenia macierzy dla $n=5$]
	$p=(4, 5, 2, 1, 2, 3)$. \\
	$A_1 \text{ - } 4 \times 5,\ A_2 \text{ - } 5 \times 2,\ A_3 \text{ - } 2 \times 1,\ A_4 \text{ - } 1 \times 2,\ A_5 \text{ - } 2 \times 3$.

	\begin{minipage}{0.45\linewidth}
		$m$
		$\begin{bNiceArray}{c|ccccc}[margin,hvlines]
				i \backslash j & 1 & 2 & 3 & 4 & 5 \\
				1 & 0 & 40 & 30 & 38 & 48 \\
				2 & \times & 0 & 10 & 20 & 31 \\
				3 & \times & \times & 0 & 4 & 12 \\
				4 & \times & \times & \times & 0 & 6 \\
				5 & \times & \times & \times & \times & 0
			\end{bNiceArray}$
	\end{minipage}\hfill
	\begin{minipage}{0.45\linewidth}
		$s$
		$\begin{bNiceArray}{c|ccccc}[margin,hvlines]
				i \backslash j & 1 & 2 & 3 & 4 & 5 \\
				1 & \times & 1 & 1 & 3 & 3 \\
				2 & \times & \times & 2 & 3 & 3 \\
				3 & \times & \times & \times & 3 & 3 \\
				4 & \times & \times & \times & \times & 4 \\
				5 & \times & \times & \times & \times & \times
			\end{bNiceArray}$
	\end{minipage}
	\vspace{1em}

	Prześledźmy, jak rosną obie tablice. Algorytm wypełnia je \emph{przekątnymi}: najpierw $L=1$ (przekątna główna), potem $L=2$, $L=3$, \dots. Komórki dopisane w~danym kroku wyróżniono na \textcolor{accentB}{pomarańczowo}; puste pola to podproblemy jeszcze nieobliczone, a~$\times$ -- pola nieużywane ($i \geq j$ dla $s$, $i > j$ dla $m$).

	\medskip
	\noindent Stan początkowy ($L=1$, pojedyncze macierze, koszt $0$):\\[0.7em]
	\begin{minipage}{0.45\linewidth}
		$m$
		$\begin{bNiceArray}{c|ccccc}[margin,hvlines]
				i \backslash j & 1 & 2 & 3 & 4 & 5 \\
				1 & 0 &  &  &  &  \\
				2 & \times & 0 &  &  &  \\
				3 & \times & \times & 0 &  &  \\
				4 & \times & \times & \times & 0 &  \\
				5 & \times & \times & \times & \times & 0
			\end{bNiceArray}$
	\end{minipage}\hfill
	\begin{minipage}{0.45\linewidth}
		$s$
		$\begin{bNiceArray}{c|ccccc}[margin,hvlines]
				i \backslash j & 1 & 2 & 3 & 4 & 5 \\
				1 & \times &  &  &  &  \\
				2 & \times & \times &  &  &  \\
				3 & \times & \times & \times &  &  \\
				4 & \times & \times & \times & \times &  \\
				5 & \times & \times & \times & \times & \times
			\end{bNiceArray}$
	\end{minipage}

	\medskip
	\noindent{\color{placeholdercolor}\rule{\linewidth}{0.4pt}}\par\smallskip
	\noindent{\color{dygresjagray}\footnotesize$\begin{array}{c|*{6}{c}} i & 0 & 1 & 2 & 3 & 4 & 5 \\ \hline p_i & 4 & 5 & 2 & 1 & 2 & 3 \end{array}$\qquad $A_i = p_{i-1}\times p_i$}
	\begin{flalign*}
		L=2: \quad & m(1, 2) = 4 \cdot 5 \cdot 2 = 40 &  & \\
		           & m(2, 3) = 5 \cdot 2 \cdot 1 = 10 &  & \\
		           & m(3, 4) = 2 \cdot 1 \cdot 2 = 4  &  & \\
		           & m(4, 5) = 1 \cdot 2 \cdot 3 = 6  &  &
	\end{flalign*}

	\medskip
	\noindent po $L=2$ (pary sąsiednich macierzy):\\[0.7em]
	\begin{minipage}{0.45\linewidth}
		$m$
		$\begin{bNiceArray}{c|ccccc}[margin,hvlines]
				i \backslash j & 1 & 2 & 3 & 4 & 5 \\
				1 & 0 & \color{accentB}40 &  &  &  \\
				2 & \times & 0 & \color{accentB}10 &  &  \\
				3 & \times & \times & 0 & \color{accentB}4 &  \\
				4 & \times & \times & \times & 0 & \color{accentB}6 \\
				5 & \times & \times & \times & \times & 0
			\end{bNiceArray}$
	\end{minipage}\hfill
	\begin{minipage}{0.45\linewidth}
		$s$
		$\begin{bNiceArray}{c|ccccc}[margin,hvlines]
				i \backslash j & 1 & 2 & 3 & 4 & 5 \\
				1 & \times & \color{accentB}1 &  &  &  \\
				2 & \times & \times & \color{accentB}2 &  &  \\
				3 & \times & \times & \times & \color{accentB}3 &  \\
				4 & \times & \times & \times & \times & \color{accentB}4 \\
				5 & \times & \times & \times & \times & \times
			\end{bNiceArray}$
	\end{minipage}

	\medskip
	\noindent{\color{placeholdercolor}\rule{\linewidth}{0.4pt}}\par\smallskip
	\noindent{\color{dygresjagray}\footnotesize$\begin{array}{c|*{6}{c}} i & 0 & 1 & 2 & 3 & 4 & 5 \\ \hline p_i & 4 & 5 & 2 & 1 & 2 & 3 \end{array}$\qquad $A_i = p_{i-1}\times p_i$}
	\begin{flalign*}
		L=3: \quad & m(1, 3) = \min \begin{cases}
			                            m(1, 1) + m(2, 3) + 4 \cdot 5 \cdot 1 \\
			                            m(1, 2) + m(3, 3) + 4 \cdot 2 \cdot 1
		                            \end{cases} = \min \begin{cases}
			                                               0 + 10 + 20 \\
			                                               40 + 0 + 8
		                                               \end{cases} = \min \begin{cases} 30 \\ 48 \end{cases} = 30 \\
		           & m(2, 4) = \min \begin{cases}
			                            m(2, 2) + m(3, 4) + 5 \cdot 2 \cdot 2 \\
			                            m(2, 3) + m(4, 4) + 5 \cdot 1 \cdot 2
		                            \end{cases} = \min \begin{cases}
			                                               0 + 4 + 20 \\
			                                               10 + 0 + 10
		                                               \end{cases} = \min \begin{cases} 24 \\ 20 \end{cases} = 20 \\
		           & m(3, 5) = \min \begin{cases}
			                            m(3, 3) + m(4, 5) + 2 \cdot 1 \cdot 3 \\
			                            m(3, 4) + m(5, 5) + 2 \cdot 2 \cdot 3
		                            \end{cases} = \min \begin{cases}
			                                               0 + 6 + 6 \\
			                                               4 + 0 + 12
		                                               \end{cases} = \min \begin{cases} 12 \\ 16 \end{cases} = 12
	\end{flalign*}

	\medskip
	\noindent po $L=3$:\\[0.7em]
	\begin{minipage}{0.45\linewidth}
		$m$
		$\begin{bNiceArray}{c|ccccc}[margin,hvlines]
				i \backslash j & 1 & 2 & 3 & 4 & 5 \\
				1 & 0 & 40 & \color{accentB}30 &  &  \\
				2 & \times & 0 & 10 & \color{accentB}20 &  \\
				3 & \times & \times & 0 & 4 & \color{accentB}12 \\
				4 & \times & \times & \times & 0 & 6 \\
				5 & \times & \times & \times & \times & 0
			\end{bNiceArray}$
	\end{minipage}\hfill
	\begin{minipage}{0.45\linewidth}
		$s$
		$\begin{bNiceArray}{c|ccccc}[margin,hvlines]
				i \backslash j & 1 & 2 & 3 & 4 & 5 \\
				1 & \times & 1 & \color{accentB}1 &  &  \\
				2 & \times & \times & 2 & \color{accentB}3 &  \\
				3 & \times & \times & \times & 3 & \color{accentB}3 \\
				4 & \times & \times & \times & \times & 4 \\
				5 & \times & \times & \times & \times & \times
			\end{bNiceArray}$
	\end{minipage}

	\medskip
	\noindent{\color{placeholdercolor}\rule{\linewidth}{0.4pt}}\par\smallskip
	\noindent{\color{dygresjagray}\footnotesize$\begin{array}{c|*{6}{c}} i & 0 & 1 & 2 & 3 & 4 & 5 \\ \hline p_i & 4 & 5 & 2 & 1 & 2 & 3 \end{array}$\qquad $A_i = p_{i-1}\times p_i$}
	\begin{flalign*}
		L=4: \quad & m(1, 4) = \min \begin{cases}
			                            m(1, 1) + m(2, 4) + 4 \cdot 5 \cdot 2 \\
			                            m(1, 2) + m(3, 4) + 4 \cdot 2 \cdot 2 \\
			                            m(1, 3) + m(4, 4) + 4 \cdot 1 \cdot 2
		                            \end{cases} = \min \begin{cases}
			                                               0 + 20 + 40 \\
			                                               40 + 4 + 16 \\
			                                               30 + 0 + 8
		                                               \end{cases} = \min \begin{cases} 60 \\ 60 \\ 38 \end{cases} = 38 \\
		           & m(2, 5) = \min \begin{cases}
			                            m(2, 2) + m(3, 5) + 5 \cdot 2 \cdot 3 \\
			                            m(2, 3) + m(4, 5) + 5 \cdot 1 \cdot 3 \\
			                            m(2, 4) + m(5, 5) + 5 \cdot 2 \cdot 3
		                            \end{cases} = \min \begin{cases}
			                                               0 + 12 + 30 \\
			                                               10 + 6 + 15 \\
			                                               20 + 0 + 30
		                                               \end{cases} = \min \begin{cases} 42 \\ 31 \\ 50 \end{cases} = 31
	\end{flalign*}

	\medskip
	\noindent po $L=4$:\\[0.7em]
	\begin{minipage}{0.45\linewidth}
		$m$
		$\begin{bNiceArray}{c|ccccc}[margin,hvlines]
				i \backslash j & 1 & 2 & 3 & 4 & 5 \\
				1 & 0 & 40 & 30 & \color{accentB}38 &  \\
				2 & \times & 0 & 10 & 20 & \color{accentB}31 \\
				3 & \times & \times & 0 & 4 & 12 \\
				4 & \times & \times & \times & 0 & 6 \\
				5 & \times & \times & \times & \times & 0
			\end{bNiceArray}$
	\end{minipage}\hfill
	\begin{minipage}{0.45\linewidth}
		$s$
		$\begin{bNiceArray}{c|ccccc}[margin,hvlines]
				i \backslash j & 1 & 2 & 3 & 4 & 5 \\
				1 & \times & 1 & 1 & \color{accentB}3 &  \\
				2 & \times & \times & 2 & 3 & \color{accentB}3 \\
				3 & \times & \times & \times & 3 & 3 \\
				4 & \times & \times & \times & \times & 4 \\
				5 & \times & \times & \times & \times & \times
			\end{bNiceArray}$
	\end{minipage}

	\medskip
	\noindent{\color{placeholdercolor}\rule{\linewidth}{0.4pt}}\par\smallskip
	\noindent{\color{dygresjagray}\footnotesize$\begin{array}{c|*{6}{c}} i & 0 & 1 & 2 & 3 & 4 & 5 \\ \hline p_i & 4 & 5 & 2 & 1 & 2 & 3 \end{array}$\qquad $A_i = p_{i-1}\times p_i$}
	\begin{flalign*}
		L=5: \quad & m(1, 5) = \min \begin{cases}
			                            m(1, 1) + m(2, 5) + 4 \cdot 5 \cdot 3 \\
			                            m(1, 2) + m(3, 5) + 4 \cdot 2 \cdot 5 \\
			                            m(1, 3) + m(4, 5) + 4 \cdot 1 \cdot 3 \\
			                            m(1, 4) + m(5, 5) + 4 \cdot 2 \cdot 3
		                            \end{cases} = \min \begin{cases}
			                                               0 + 31 + 60  \\
			                                               40 + 12 + 24 \\
			                                               30 + 6 + 12  \\
			                                               38 + 0 + 24
		                                               \end{cases} = \min \begin{cases} 91 \\ 76 \\ 48 \\ 62 \end{cases} = 48
	\end{flalign*}

	\medskip
	\noindent po $L=5$ (cały ciąg -- ostatnia komórka):\\[0.7em]
	\begin{minipage}{0.45\linewidth}
		$m$
		$\begin{bNiceArray}{c|ccccc}[margin,hvlines]
				i \backslash j & 1 & 2 & 3 & 4 & 5 \\
				1 & 0 & 40 & 30 & 38 & \color{accentB}48 \\
				2 & \times & 0 & 10 & 20 & 31 \\
				3 & \times & \times & 0 & 4 & 12 \\
				4 & \times & \times & \times & 0 & 6 \\
				5 & \times & \times & \times & \times & 0
			\end{bNiceArray}$
	\end{minipage}\hfill
	\begin{minipage}{0.45\linewidth}
		$s$
		$\begin{bNiceArray}{c|ccccc}[margin,hvlines]
				i \backslash j & 1 & 2 & 3 & 4 & 5 \\
				1 & \times & 1 & 1 & 3 & \color{accentB}3 \\
				2 & \times & \times & 2 & 3 & 3 \\
				3 & \times & \times & \times & 3 & 3 \\
				4 & \times & \times & \times & \times & 4 \\
				5 & \times & \times & \times & \times & \times
			\end{bNiceArray}$
	\end{minipage}

	\medskip
	Każdy przebieg pętli po~$l$ dopisuje dokładnie jedną przekątną: w~komórce $m(i, j)$ ląduje koszt, a~w~bliźniaczej $s(i, j)$ -- zwycięskie cięcie $k$. Ostatnią obliczoną wartością jest $m(1, 5) = 48$ w~prawym górnym rogu -- szukany koszt całego ciągu.

	Oto algorytm znajdowania optymalnego rozwiązania:
	\begin{alg}{Algorytm Matrix Chain Multiply}{alg:matrixchainmultiply}
		\FunctionNN{MatrixChainMultiply}{$A, s, i, j$} \Comment{$A = (A_1, \dots, A_n)$}
		\If{$j = i$}
		\State \Return $A_i$
		\EndIf
		\State $X$ \Assign \Call{MatrixChainMultiply}{$A, s, i, s(i, j)$}
		\State $Y$ \Assign \Call{MatrixChainMultiply}{$A, s, s(i, j) + 1, j$}
		\State \Return \Call{MatrixMultiply}{$X, Y$}
		\EndFunction
	\end{alg}

	W celu znalezienia rozwiązania wywołujemy: \Call{MatrixChainMultiply}{$A, s, 1, n$}
	\begin{flalign*}
		 & A_1 \cdot A_2 \cdot A_3 \cdot A_4 \cdot A_5                                                                  &  & \\
		 & s(1, 5) = 3 \implies (A_1 \cdot A_2 \cdot A_3)(A_4 \cdot A_5)                                                &  & \\
		 & s(1, 3) = 1 \implies (A_1 \cdot (A_2 \cdot A_3))(A_4 \cdot A_5) \qquad s(4, 5) = 4 \implies \text{bez zmian} &  & \\
		 & s(2, 3) = 2 \implies (A_1 \cdot (A_2 \cdot A_3))(A_4 \cdot A_5)                                              &  &
	\end{flalign*}
\end{example}

Aby algorytm zadziałał dobrze, problem musi spełniać warunki:
\begin{enumerate}
	\item \uemph{Własność optymalnej podstruktury} -- optymalne rozwiązanie problemu można skonstruować z optymalnych rozwiązań podproblemów.
	\item \uemph{Własność wspólnych podproblemów} -- sytuacja, gdy prosty algorytm rekurencyjny obliczałby rozwiązania tych samych podproblemów wielokrotnie. Liczba podproblemów powinna być ,,mała''.
\end{enumerate}

W klasycznej metodzie programowania dynamicznego rozwiązanie jest budowane od dołu do góry -- od podproblemów najmniejszych do największych. Jest też podejście, w którym rozwiązanie jest budowane od góry do dołu, podobnie jak w prostym algorytmie rekurencyjnym, ale z tą różnicą, że unikamy obliczania tych samych podproblemów wielokrotnie. To podejście nazywa się \uemph{spamiętywaniem} (\textit{memoization}).

\paragraph{Problem mnożenia ciągu macierzy ze spamiętywaniem}
Algorytm oblicza tablice $m$ i $s$ jak poprzednio, ale w inny sposób. Dopóki podproblem nie jest rozwiązany, $m(i, j)$ przechowuje $\infty$. Po rozwiązaniu podproblemu $(i, j)$ umieszczamy w $m(i, j)$ odpowiednią wartość i od tej pory używamy jej bez obliczania.

\begin{alg}{Algorytm Lookup Chain}{alg:lookupchain}
	\FunctionNN{LookupChain}{$p, i, j$}
	\Comment{$p = (p_0, \dots, p_n)$ - ciąg wymiarów; $A_i$ - macierz $p_{i-1} \times p_i$}
	\If{$m(i, j) < \infty$}
	\State \Return $m(i, j), s(i, j)$
	\EndIf
	\If{$i = j$}
	\State $m(i, j) \Assign 0$
	\Else
	\For{$k \Assign i$ \To $j - 1$}
	\State $q \Assign \Call{LookupChain}{p, i, k} + \Call{LookupChain}{p, k + 1, j} + p_{i-1} \cdot p_k \cdot p_j$
	\If{$q < m(i, j)$}
	\State $m(i, j) \Assign q$
	\State $s(i, j) \Assign k$
	\EndIf
	\EndFor
	\EndIf
	\State \Return $(m(i, j), s(i, j))$
	\EndFunction
\end{alg}

\begin{alg}{Algorytm Memoized Matrix Chain}{alg:memoizedmatrixchain}
	\FunctionNN{MemoizedMatrixChain}{$p$}
	\State $n \Assign \Length{p} - 1$
	\For{$i \Assign 1$ \To $n$}
	\For{$j \Assign i$ \To $n$}
	\State $m(i, j) \Assign \infty$
	\EndFor
	\EndFor
	\State \Return \Call{LookupChain}{$p, 1, n$}
	\EndFunction
\end{alg}

\paragraph{Złożoność czasowa:}
\begin{itemize}
	\item Linie 1--4 w \Call{MemoizedMatrixChain}{} zajmują $\BT(n^2)$.
	\item Każdy z $\BO(n^2)$ elementów tablicy $m$ jest modyfikowany raz, później odczytanie wartości $m(i, j)$ zajmuje czas stały.
	\item Modyfikacja $m(i, j)$ zajmuje $\BO(n)$ czasu (pętla w liniach 5--9 w \Call{LookupChain}{}).
\end{itemize}
Ostatecznie złożoność algorytmu wynosi $\BO(n^3)$.


% ════════════════════════════════════════════════════════════════
\subsection{Problem najdłuższego wspólnego podciągu}
% ════════════════════════════════════════════════════════════════

Niech $X = (x_1, \dots, x_m)$ -- ciąg. \\
$Z = (z_1, \dots, z_k)$ jest podciągiem $X$, jeśli istnieje rosnący ciąg indeksów $1 \leq i_1 < i_2 < \dots < i_k \leq m$ takich, że dla każdego $j \in \{1, \dots, k\}$, $z_j = x_{i_j}$.

\begin{example}[Podciąg]
	$X = (A, {\color{accent}B}, {\color{accent}C}, B, {\color{accent}D}, A, {\color{accent}B})$, $Z =
		({\color{accent}B}, {\color{accent}C}, {\color{accent}D}, {\color{accent}B})$.
	Indeksy: $(2, 3, 5, 7)$.
\end{example}

\begin{problem}[najdłuższego wspólnego podciągu]
	Dane są ciągi $X = (x_1, \dots, x_m)$ oraz $Y = (y_1, \dots, y_n)$.
	Należy znaleźć najdłuższy wspólny podciąg (LCS -- \textit{Longest Common Subsequence}) $X$ i $Y$, czyli najdłuższy ciąg, który jest podciągiem zarówno $X$, jak i $Y$.
\end{problem}

Wszystkich podciągów $X$ jest $2^m$, więc algorytm sprawdzający wszystkie podciągi $X$ działa w czasie wykładniczym.

Niech $X_i = (x_1, \dots, x_i)$ dla $i \in \{1, \dots, m\}$ i dodatkowo $X_0 = ()$ (ciąg pusty). Analogicznie dla $Y$.

\begin{lemma}
	Niech $X = (x_1, \dots, x_m)$, $Y = (y_1, \dots, y_n)$ i niech $Z = (z_1, \dots, z_k)$ będzie LCS $X$ i $Y$.
	\begin{enumerate}
		\item Jeśli $x_m = y_n$, to $z_k = x_m = y_n$ i $Z_{k-1}$ jest LCS $X_{m-1}$ i $Y_{n-1}$.
		\item Jeśli $x_m \neq y_n$ i $z_k \neq x_m$, to $Z$ jest LCS $X_{m-1}$ i $Y$.
		\item Jeśli $x_m \neq y_n$ i $z_k \neq y_n$, to $Z$ jest LCS $X$ i $Y_{n-1}$.
	\end{enumerate}
\end{lemma}

\tikzset{
	lcs/.style={font=\small, baseline=(current bounding box.center)},
	lab/.style={anchor=east, font=\itshape},
	pref/.style={draw, rounded corners=2pt, fill=accent!8, minimum width=1.7cm, minimum height=0.62cm, anchor=west, inner sep=2pt},
	cell/.style={draw, minimum width=0.62cm, minimum height=0.62cm, anchor=west, inner sep=1pt},
	match/.style={fill=accentA!55},
	drop/.style={fill=placeholdercolor!25, text=dygresjagray},
	sub/.style={accentB, thick, dashed, rounded corners=2pt},
}

\begin{figure}[ht]
	\centering
	% ----- Przypadek 1: x_m = y_n -----
	\begin{minipage}[t]{0.31\textwidth}
		\centering
		\begin{tikzpicture}[lcs]
			\node[lab] at (-0.4,0) {$X$};
			\node[pref] (Xp) at (0,0) {$X_{m-1}$};
			\node[cell, match] (xm) at (Xp.east) {$x_m$};
			\node[lab] at (-0.4,-0.95) {$Y$};
			\node[pref] (Yp) at (0,-0.95) {$Y_{n-1}$};
			\node[cell, match] (yn) at (Yp.east) {$y_n$};
			\draw[accentB, thick] (xm.south) -- (yn.north)
			node[midway, fill=white, inner sep=1pt, text=accentB] {$=$};
			\draw[sub] ([shift={(-2pt,2pt)}]Xp.north west) rectangle ([shift={(2pt,-2pt)}]Xp.south east);
			\draw[sub] ([shift={(-2pt,2pt)}]Yp.north west) rectangle ([shift={(2pt,-2pt)}]Yp.south east);
		\end{tikzpicture}
		\par\smallskip
		{\footnotesize $x_m = y_n$\par
			$z_k = x_m = y_n$ wchodzi do LCS; \\ podproblem $(X_{m-1}, Y_{n-1})$}
	\end{minipage}\hfill
	% ----- Przypadek 2: x_m != y_n, z_k != x_m -----
	\begin{minipage}[t]{0.31\textwidth}
		\centering
		\begin{tikzpicture}[lcs]
			\node[lab] at (-0.4,0) {$X$};
			\node[pref] (Xp) at (0,0) {$X_{m-1}$};
			\node[cell, drop] (xm) at (Xp.east) {$x_m$};
			\draw[dygresjagray] (xm.south west) -- (xm.north east);
			\node[lab] at (-0.4,-0.95) {$Y$};
			\node[pref] (Yp) at (0,-0.95) {$Y_{n-1}$};
			\node[cell] (yn) at (Yp.east) {$y_n$};
			\draw[sub] ([shift={(-2pt,2pt)}]Xp.north west) rectangle ([shift={(2pt,-2pt)}]Xp.south east);
			\draw[sub] ([shift={(-2pt,2pt)}]Yp.north west) rectangle ([shift={(2pt,-2pt)}]yn.south east);
		\end{tikzpicture}
		\par\smallskip
		{\footnotesize $x_m \neq y_n$ i $z_k \neq x_m$\par
			$x_m$ nie należy do LCS; \\ podproblem $(X_{m-1}, Y)$}
	\end{minipage}\hfill
	% ----- Przypadek 3: x_m != y_n, z_k != y_n -----
	\begin{minipage}[t]{0.31\textwidth}
		\centering
		\begin{tikzpicture}[lcs]
			\node[lab] at (-0.4,0) {$X$};
			\node[pref] (Xp) at (0,0) {$X_{m-1}$};
			\node[cell] (xm) at (Xp.east) {$x_m$};
			\node[lab] at (-0.4,-0.95) {$Y$};
			\node[pref] (Yp) at (0,-0.95) {$Y_{n-1}$};
			\node[cell, drop] (yn) at (Yp.east) {$y_n$};
			\draw[dygresjagray] (yn.south west) -- (yn.north east);
			\draw[sub] ([shift={(-2pt,2pt)}]Xp.north west) rectangle ([shift={(2pt,-2pt)}]xm.south east);
			\draw[sub] ([shift={(-2pt,2pt)}]Yp.north west) rectangle ([shift={(2pt,-2pt)}]Yp.south east);
		\end{tikzpicture}
		\par\smallskip
		{\footnotesize $x_m \neq y_n$ i $z_k \neq y_n$\par
			$y_n$ nie należy do LCS; \\ podproblem $(X, Y_{n-1})$}
	\end{minipage}
	\caption{Trzy przypadki lematu o~optymalnej podstrukturze LCS. Żółta komórka --
		dopasowany ostatni element ($x_m = y_n$), który wchodzi do LCS; szara,
		przekreślona komórka -- element odrzucony, nienależący do tego LCS;
		przerywana pomarańczowa ramka -- podproblem(y), do których redukuje się zadanie.}
	\label{fig:lcs:cases}
\end{figure}

Rozwiązanie dla problemu $(X, Y)$ można skonstruować z rozwiązań podproblemów $(X_{m-1}, Y_{n-1})$, $(X_{m-1}, Y)$ i $(X, Y_{n-1})$ -- \textbf{własność optymalnej podstruktury}.
Podproblem $(X_{m-1}, Y_{n-1})$ jest zawarty w podproblemach $(X_{m-1}, Y)$ i $(X, Y_{n-1})$ -- \textbf{własność wspólnych podproblemów}.

Niech $c(i, j)$ oznacza długość LCS ciągów $X_i$ i $Y_j$.
\begin{flalign*}
	c(i, j) = \begin{cases}
		          0                                & \text{jeśli } i = 0 \text{ lub } j = 0         \\
		          c(i - 1, j - 1) + 1              & \text{jeśli } i, j > 0 \text{ i } x_i = y_j    \\
		          \max\{c(i - 1, j), c(i, j - 1)\} & \text{jeśli } i, j > 0 \text{ i } x_i \neq y_j
	          \end{cases} &  &  &  &
\end{flalign*}

Będziemy też używać tablicy $b(i, j)$, aby znaleźć sam podciąg LCS.
Tablice mają wymiary: $c[0 \dots m, 0 \dots n]$ oraz $b[1 \dots m, 1 \dots n]$. Wartości w tablicy kierunków to $b(i, j) \in \{\nwarrow, \uparrow, \leftarrow\}$.

\begin{alg}{Algorytm obliczający długość LCS}{alg:lcs}
	\FunctionNN{LCS}{$X, Y$}
	\State $m \Assign \Length{X}$
	\State $n \Assign \Length{Y}$
	\For{$i \Assign 0$ \To $m$}
	\State $c(i, 0) \Assign 0$
	\EndFor
	\For{$j \Assign 0$ \To $n$}
	\State $c(0, j) \Assign 0$
	\EndFor
	\For{$i \Assign 1$ \To $m$}
	\For{$j \Assign 1$ \To $n$}
	\If{$x_i = y_j$}
	\State $c(i, j) \Assign c(i - 1, j - 1) + 1$
	\State $b(i, j) \Assign \text{,,}\nwarrow\text{''}$
	\ElsIf{$c(i - 1, j) \geq c(i, j - 1)$}
	\State $c(i, j) \Assign c(i - 1, j)$
	\State $b(i, j) \Assign \text{,,}\uparrow\text{''}$
	\Else
	\State $c(i, j) \Assign c(i, j - 1)$
	\State $b(i, j) \Assign \text{,,}\leftarrow\text{''}$
	\EndIf
	\EndFor
	\EndFor
	\State \Return $(c, b)$
	\EndFunction
\end{alg}

\paragraph{Przykład:}
$X = (A, {\color{accent}B}, {\color{accent}C}, {\color{accent}B}, D, {\color{accent}A}, B)$, $Y =
	({\color{accent}B}, D, {\color{accent}C}, A, {\color{accent}B}, {\color{accent}A})$. \\
$m = 7$, $n = 6$.

\begin{table}[H]
	\centering
	\renewcommand{\arraystretch}{1.2}
	\begin{minipage}{0.48\linewidth}
		\centering
		Tablica kosztów $c$: \\[0.5em]
		\begin{NiceTabular}{cc|c|cccccc}[margin,hvlines]
			    & $j$   & 0     & 1   & 2   & 3   & 4   & 5   & 6   \\
			$i$ &       & $y_j$ & $B$ & $D$ & $C$ & $A$ & $B$ & $A$ \\
			\hline
			0   & $x_i$ & 0     & 0   & 0   & 0   & 0   & 0   & 0   \\
			1   & $A$   & 0     & 0   & 0   & 0   & 1   & 1   & 1   \\
			2   & $B$   & 0     & 1   & 1   & 1   & 1   & 2   & 2   \\
			3   & $C$   & 0     & 1   & 1   & 2   & 2   & 2   & 2   \\
			4   & $B$   & 0     & 1   & 1   & 2   & 2   & 3   & 3   \\
			5   & $D$   & 0     & 1   & 2   & 2   & 2   & 3   & 3   \\
			6   & $A$   & 0     & 1   & 2   & 2   & 3   & 3   & 4   \\
			7   & $B$   & 0     & 1   & 2   & 2   & 3   & 4   & 4   \\
		\end{NiceTabular}
	\end{minipage}\hfill
	\begin{minipage}{0.48\linewidth}
		\centering
		Tablica strzałek $b$: \\[0.5em]
		\begin{NiceTabular}{cc|c|cccccc}[margin,hvlines]
			    & $j$   & 0     & 1          & 2            & 3            & 4            & 5            & 6            \\
			$i$ &       & $y_j$ & $B$        & $D$          & $C$          & $A$          & $B$          & $A$          \\
			\hline
			0   & $x_i$ &       &            &              &              &              &              &              \\
			1   & $A$   &       & $\uparrow$ & $\uparrow$   & $\uparrow$   & $\nwarrow$   & $\leftarrow$ & $\nwarrow$   \\
			2   & $B$   &       & $\nwarrow$ & $\leftarrow$ & $\leftarrow$ & $\uparrow$   & $\nwarrow$   & $\leftarrow$ \\
			3   & $C$   &       & $\uparrow$ & $\uparrow$   & $\nwarrow$   & $\leftarrow$ & $\uparrow$   & $\uparrow$   \\
			4   & $B$   &       & $\nwarrow$ & $\uparrow$   & $\uparrow$   & $\uparrow$   & $\nwarrow$   & $\leftarrow$ \\
			5   & $D$   &       & $\uparrow$ & $\nwarrow$   & $\uparrow$   & $\uparrow$   & $\uparrow$   & $\uparrow$   \\
			6   & $A$   &       & $\uparrow$ & $\uparrow$   & $\uparrow$   & $\nwarrow$   & $\uparrow$   & $\nwarrow$   \\
			7   & $B$   &       & $\nwarrow$ & $\uparrow$   & $\uparrow$   & $\uparrow$   & $\nwarrow$   & $\uparrow$   \\
		\end{NiceTabular}
	\end{minipage}
\end{table}

\begin{alg}{Wypisywanie najdłuższego wspólnego podciągu}{alg:printlcs}
	\FunctionNN{PrintLCS}{$b, X, Y, i, j$}
	\If{$i = 0$ \textbf{lub} $j = 0$}
	\State \Return $\varepsilon$ \Comment{ciąg pusty}
	\EndIf
	\If{$b(i, j) = \text{,,}\nwarrow\text{''}$}
	\State \Call{PrintLCS}{$b, X, Y, i - 1, j - 1$}
	\State \textbf{print} $x_i$
	\ElsIf{$b(i, j) = \text{,,}\uparrow\text{''}$}
	\State \Call{PrintLCS}{$b, X, Y, i - 1, j$}
	\Else
	\State \Call{PrintLCS}{$b, X, Y, i, j - 1$}
	\EndIf
	\EndFunction
\end{alg}

Żeby wypisać rozwiązanie, wywołujemy \Call{PrintLCS}{$b, X, Y, m, n$}. \\
Dla powyższego przykładu otrzymamy: $(B, C, B, A)$.

Złożoność czasowa tego algorytmu wynosi $\BT(m \cdot n)$.
